% !TeX program = pdflatex
%%% doc/bangor.tex — documentation for the bangor package family
\documentclass[a4paper,12pt]{article}
\usepackage{bangor}
\usepackage{bangoridentity}
\usepackage{bangorlayout}
\usepackage{bangordeclarations}
\usepackage[hidelinks]{hyperref}
\usepackage{enumitem}

\title{The bangor package family}
\author{Bangor University LaTeX templates project}
\date{2026/09/19}
\hypersetup{pdfsubject={\bangorversionstamp},
  pdfkeywords={Bangor University; \bangorversionstamp}}

\begin{document}
\maketitle

\section{What this package does}
The \textsf{bangor} family typesets documents for Bangor University.
It provides the brand, the document identity, the regulation layout,
the statutory declarations, a headed letter style, and a thesis class.
This document is written in controlled English, so that readers who
do not use English as a first language can read it without difficulty.

\section{The layers}
The family has three layers. A module must only use modules in a lower
layer. Do not make a cycle between modules.

\begin{enumerate}
\item \textsf{bangor} holds the brand facts: the colours, the crest,
and the wordmark. It holds no document logic.
\item \textsf{bangoridentity}, \textsf{bangorlayout}, and
\textsf{bangordeclarations} are reusable document modules. You can use
each one with any base class.
\item \textsf{bangorletter} and the \textsf{bangorthesis} class are
consumers. They use the lower layers and add their own document type.
\end{enumerate}

\section{The brand layer: bangor}
Load the package with \verb|\usepackage{bangor}|. It gives you three
things.

\begin{itemize}
\item The colours \verb|bangorred| and \verb|bangorsun|, from the
university brand hub.
\item The crest: \verb|\bangorcrest[3cm]|. The optional argument sets
the width. The default width is 3cm.
\item The wordmark: \verb|\bangorwordmark|. This prints the university
name in the brand colour and weight.
\end{itemize}

The version stamp macro \verb|\bangorversionstamp| lists every loaded
module and its version. Consumers place this stamp in the PDF metadata,
so every generated PDF records what produced it.

\section{Document identity: bangoridentity}
Set the document facts in the preamble. School and college are free
text, because the university renames its schools and colleges.

\begin{itemize}
\item \verb|\school{School of Computer Science and Engineering}|
\item \verb|\college{College of Science and Engineering}|
\item \verb|\degreeScheme{Doctor of Philosophy}|
\item \verb|\supervisor{Dr A. Supervisor}| --- call this command once
for each supervisor.
\item \verb|\submissiondate{September 2026}|
\item \verb|\wordcount{9930}| --- for taught work that must state the
word count.
\end{itemize}

The command \verb|\bangoridentityblock| renders the crest, the
wordmark, and the school and college lines. The thesis title page and
the letter letterhead both use this block, so neither repeats the brand
facts.

\section{Regulation layout: bangorlayout}
This module applies the rules from Regulation 03, 2025 Version 01,
section 6.3, and checks them.

\begin{itemize}
\item Fonts: \verb|font=termes| (Times-Roman equivalent, the default),
\verb|font=bonum| (Bookman equivalent), or \verb|font=heros| (Arial
and Helvetica equivalent, sans-serif). Serif body text must be at least
12pt. Sans-serif body text must be at least 10pt.
\item Margins: 2.5cm at the binding edge, 2cm on the other sides.
\item Line spacing: 1.5 in the main text. Set the abstract, indented
quotations, and footnotes single spaced if you want; the regulation
permits this.
\end{itemize}

Strict mode is on by default. A rule that the module does not meet
fails the build. Set \verb|strict=false| to turn each failure into a
warning. The module checks the settings that it controls. It cannot
check changes that you make around it.

\section{Statutory declarations: bangordeclarations}
The command \verb|\statementspage| prints the Statement of Originality
and the Statement of Availability. The originality text reproduces the
statutory declaration from the regulation exactly. Load the package
with the \verb|welsh| option to print the Welsh text as well.

The generative AI statement is a fill-in template. Replace it with your
own statement:

\begin{quote}
\verb|\aistatement{Your statement about your use of AI tools.}|
\end{quote}

You can use this package with any base class that provides
\verb|\section*|, without the thesis class.

\section{Headed letters: bangorletter}
Use the standard \textsf{letter} class, then load
\verb|\usepackage{bangorletter}|. Put \verb|\bangorletterhead| at the
top of each letter, before \verb|\opening|. The command prints the
crest, the wordmark, the school line if you set one, and a rule in the
brand red.

\section{The thesis class: bangorthesis}
Start with \verb|\documentclass[phd,12pt,a4paper]{bangorthesis}|.

\subsection{Class options}
\begin{itemize}
\item Degree keywords: \verb|phd|, \verb|mphil|, \verb|mres|,
\verb|mscres|, \verb|engd|, and \verb|profdoc| set the wording
``thesis''. The taught keywords, for example \verb|msc|, \verb|bsc|,
and \verb|beng|, set the wording ``dissertation''. The options
\verb|thesis| and \verb|dissertation| override the wording directly.
\item \verb|citations=style| selects the biblatex style. The default
is \verb|ieee|. Any biblatex style name is valid. Ask your supervisor
which style your school requires.
\item \verb|font=termes|, \verb|strict|, and \verb|strict=false| pass
through to the layout module.
\item \verb|print| adds a blank page before the title page, for duplex
printing. \verb|digital| is the default and prints no blank page.
\item \verb|markdraft| prints DRAFT in the running header.
\item \verb|welsh| prints the statutory declaration in Welsh and loads
babel with Welsh.
\item All other options pass to the \textsf{report} class.
\end{itemize}

\subsection{Commands}
Set the metadata from \textsf{bangoridentity} in the preamble, then
call \verb|\bibliographySetup| once, add your resources with
\verb|\addbibresource|, and build the front matter:

\begin{quote}
\verb|\maketitle| \quad \verb|\statementspage| \quad
\verb|\begin{acknowledgements}| \verb|...| \verb|\end{acknowledgements}|
\quad \verb|\begin{abstract}| \verb|...| \verb|\end{abstract}| \quad
\verb|\tables|
\end{quote}

The abstract is single spaced. The regulation caps it at 600 words.

\section{Version stamping in generated PDFs}
Every consumer places the version stamp in the PDF metadata, in the
subject and the keywords fields. The stamp lists each loaded module
and its version, and the template version when the release process
provides one. Open the document information panel of the PDF to see
which versions produced it.

\end{document}
